\section{Experiments}
\label{sec:results}

Each of 32 seeded streams contains 16 independent boards. Twelve receive
the upgrade and four keep their old solver. Each board has a uniform random
goal and residue, and a random row order fixed across policies. We label
every state crossed with $+1,+3,-2,+25$, including the goal state.
The always-illegal $+25$ gives an explicit invalid family in addition to
out-of-bounds cases of the other moves. The primary release adds $+6$;
the secondary release instead adds $-5$ to the same paired boards.
Each cap is applied independently per upgraded board. A guard lookup,
rerun of a legality check, or attempted continuation move costs one
operation; no method can charge a partially finished row as refreshed.
The outcome is exact joint labels, not reward accuracy on solver outputs.

\paragraph{Positive control.}
The independent labeler expands reached-state sets forward, rather than
memoizing backward reachability. All three refresh arms match it on all
1,536 arm--board checks for each release when the cap is removed, including
unreleased boards. Freezing the store leaves 1,429 of 36,864 upgraded rows
wrong after release A (3.88\%) and 2,497 wrong after release B (6.77\%).
The primary release therefore exceeds the required 1-point stale-label
error before testing any proposed improvement.

\begin{table}[h]
\centering
\caption{\textbf{Bypassing the guard uses more of the cap for continuation work.}
Release A, primary cap of 96 operations per upgraded board. Exact joint
accuracy is a mean with 95\% marginal $t$ intervals over 32 streams; work
columns are means per upgraded board. Paired contrasts are given in text.}
\label{tab:primary}
\begin{tabular}{lrrrrr}
\toprule
Policy & Exact (\%) $\uparrow$ & Total & Guard & Validity & Continuation \\
\midrule
\csname @@input\endcsname table_accuracy.tex
\bottomrule
\end{tabular}
\end{table}

\paragraph{Primary test.}
At 96 operations per upgraded board, guarded split refresh gains
0.019 percentage points over whole-label refresh (97.5\% paired
interval: 0.001 to 0.037), but loses 0.366 points to memoized PAV
(\(-0.492\) to \(-0.241\)). These are the preregistered two-sided
intervals, giving 95\% familywise coverage by Bonferroni. Both upper
bounds fall below the required 2-point gain. This is not an inability to
resolve a noisy tie: the interval against memoized PAV lies entirely below
zero, and even the upper bound against whole-label refresh is close to zero
relative to the decision threshold.

\paragraph{Retention.}
At the primary cap, all four policies retained all 12,288 unchanged rows,
all 15,360 invalid rows across both upgraded and unchanged boards, and all
3,072 unchanged statewise weak orderings, including ties.
All 986 strict detours (unchanged rows with $P_1=-1$) retained their old
label exactly. These are intended invariances: the release does not touch task
legality or the solver on those boards. They check for collateral damage to
the label store, not generalization of a learned verifier.

\subsection{Where did the budget go?}
\label{sec:diagnostic}

Both guards mark 9,298 rows as potentially stale, 36.7\% of the legal
rows on upgraded boards. They cover all actual changes, but only 15.4\%
of flags correspond to changed labels in release A and 26.9\% in B.
The guard detects access to a new move rather than a change in reachability,
and progress can stay the same even when a reachability term changes.
Guard charges at the primary cap cost split refresh 8.76 operations per
board on average, while whole-label refresh spent just 2.24 on legality.
The gain from skipping that check is small in an arithmetic grammar; the
unguarded baseline spends almost all available work on shared continuation
queries instead (\cref{tab:primary}).

At the 384-operation cap in release A, guarded split refresh and memoized
recomputation both recover every joint label and spend the same 180.18
mean continuation attempts. The guarded version spends 66 more operations
per board, exactly one lookup per legal row. This does not imply pruning
can never save search. In this store, unflagged rows did not add to the
aggregate continuation work once queries were shared, so pruning those
rows did not save enough work to amortize the guard.

\paragraph{Second release.}
The negative move creates more stale labels but does not reverse the
comparison at the primary cap: split leads whole-label refresh by 0.016
points and trails memoized recomputation by 0.060 points. At the largest
cap split and memoized recomputation reach exact labels while whole
remains at 99.82\%. These comparisons are descriptive, not a second
promotion test. \Cref{fig:cap} shows the complete budget sweep.
